% !TeX root = ../master.tex

\chapter{Discussion}
\label{ch:discussion}

\section{Integration of Results}
\label{sec:integration}

% - Synthesize findings across all three parts
% - RQ1 answer: do LLMs rely on argument structure? (expected: yes, large delta)
% - RQ2 answer: does context help? (expected: yes, especially guidelines)
% - RQ3 answer: does fine-tuning reintroduce shortcuts? (open question)
% - The big picture: is shortcut learning architecture-dependent or training-dependent?

\section{Comparison to Encoder Baselines}
\label{sec:encoder_comparison}

% - Systematic comparison table: encoders (Feger et al.) vs. LLMs (this thesis)
% - Dimensions: absolute F1, delta under manipulation, cross-dataset transfer
% - Architectural explanation: causal attention vs. bidirectional attention
% - Do LLMs actually "understand" arguments, or just different shortcuts?

\section{The Role of Context}
\label{sec:role_of_context}

% - When does context help vs. hurt?
% - Dataset-specific definitions vs. generic: the annotation divergence hypothesis
% - Guidelines as a form of task specification vs. additional noise
% - Document context: disambiguation value vs. prompt length overhead
% - Implications for GAIC task design and future shared tasks

\section{Implications for Argument Mining}
\label{sec:implications}

% - Practical recommendations: which model + context combination for production?
% - Zero-shot LLMs as a viable alternative to fine-tuned encoders?
% - Cost-performance tradeoffs (API costs, latency, reproducibility)
% - Guidelines-as-prompts as a general strategy for task transfer

\section{Limitations}
\label{sec:limitations}

% - Sample size constraints (n=30 per dataset in many experiments)
% - Zero-shot vs. fine-tuned comparison is not perfectly controlled
%   (different absolute performance levels)
% - Contamination risk: LLMs may have seen training data
% - Only binary argument identification, not full AM pipeline
% - Prompt sensitivity: results may vary with different prompt formulations
% - Context extraction via GPT introduces another model's biases
% - English-only evaluation
